% GENERATED FILE. Edit canonical lessons, then run npm run generate:books.
% canonical-source-hash: fnv1a64:d45a8fc545fe6e72
% canonical-lessons: GU-C166-samarvu, GU-C166-cholvu, GU-C166-dungli, GU-C166-kothmir, GU-C166-vatvu

\chapter{Chopping, peeling and grinding}
\label{ch:gu-chopping-peeling-and-grinding}
\hlchaptermodality{166}

\begin{chapteropening}
\textbf{By the end of this chapter:} \emph{I can describe chopping, peeling and grinding food.}

Complete the last lesson of chapter 166: I can describe chopping, peeling and grinding food.
\end{chapteropening}

\section[\texorpdfstring{\gu{સમારવું} (samārvũ)}{સમારવું (samārvũ)}]{\gu{સમારવું} (samārvũ) — to chop, to cut up}

\label{lesson:GU-C166-samarvu}



\begin{quote}
Before the new one: say the Gujarati for to mix, then the Gujarati for to boil.
\end{quote}

\subsection*{You'll want to know: \gu{સમારવું}}
\textbf{\gu{સમારવું}} — \emph{samārvũ} — \textquotedblleft{}to chop, to cut up\textquotedblright{}.

\begin{quote}
\textbf{\gu{તે} \gu{છરીથી} \gu{શાક} \gu{સમારે} \gu{છે}.} — \emph{te charīthī śāk samāre che} — \textquotedblleft{}She chops the vegetables with a knife.\textquotedblright{}
\end{quote}

\begin{cousinweb}
\textbf{\gu{સમારવું}} (\emph{samārvũ}) is a tadbhava. CDIAL lists it, glossed \textquotedblleft{}to repair, adorn; prepare (vegetables) for cooking\textquotedblright{}, under Sanskrit \emph{saṁvārayati}, \textquotedblleft{}keeps off; puts in order\textquotedblright{} (CDIAL 13021).
\end{cousinweb}

\begin{grammarlens}[title={What you've built}]
One more word for everyday life at A2.
\end{grammarlens}

\subsection*{Your turn}
\begin{itemize}
\raggedright
  \item \emph{Say it:} \emph{samārvũ}
  \item \emph{Say it:} \emph{samārvũ}, once more
  \item \emph{Say it:} \emph{te charīthī śāk samāre che} — She chops the vegetables with a knife.
  \item \emph{Recall:} say the Gujarati for to mix, then the Gujarati for to boil, then say \emph{samārvũ} again
\end{itemize}

\begin{tcolorbox}[breakable,colback=teal!4,colframe=teal!35!black,title={Before you move on}]
What does \gu{સમારવું} mean? (\textquotedblleft{}To chop, to cut up\textquotedblright{}.) Say it once more.
\end{tcolorbox}
\section[\texorpdfstring{\gu{છોલવું} (cholvũ)}{છોલવું (cholvũ)}]{\gu{છોલવું} (cholvũ) — to peel}

\label{lesson:GU-C166-cholvu}



\begin{quote}
Before the new one: say the Gujarati for to boil, then the Gujarati for to chop.
\end{quote}

\subsection*{You'll want to know: \gu{છોલવું}}
\textbf{\gu{છોલવું}} — \emph{cholvũ} — \textquotedblleft{}to peel\textquotedblright{}.

\begin{quote}
\textbf{\gu{હું} \gu{બટાકા} \gu{છોલું} \gu{છું}.} — \emph{hũ baṭākā cholũ chũ} — \textquotedblleft{}I peel the potatoes.\textquotedblright{}
\end{quote}

\begin{cousinweb}
\textbf{\gu{છોલવું}} (\emph{cholvũ}) is a tadbhava. It comes from \emph{*chōll}, \textquotedblleft{}to peel, skin\textquotedblright{}, a reconstructed form that Burrow, cited in CDIAL, traces to a Dravidian source; CDIAL also compares Sanskrit \emph{challi}, \textquotedblleft{}bark\textquotedblright{} (CDIAL 5073).
\end{cousinweb}

\begin{grammarlens}[title={What you've built}]
One more word for everyday life at A2.
\end{grammarlens}

\subsection*{Your turn}
\begin{itemize}
\raggedright
  \item \emph{Say it:} \emph{cholvũ}
  \item \emph{Say it:} \emph{cholvũ}, once more
  \item \emph{Say it:} \emph{hũ baṭākā cholũ chũ} — I peel the potatoes.
  \item \emph{Recall:} say the Gujarati for to boil, then the Gujarati for to chop, then say \emph{cholvũ} again
\end{itemize}

\begin{tcolorbox}[breakable,colback=teal!4,colframe=teal!35!black,title={Before you move on}]
What does \gu{છોલવું} mean? (\textquotedblleft{}To peel\textquotedblright{}.) Say it once more.
\end{tcolorbox}
\section[\texorpdfstring{\gu{ડુંગળી} (ḍungḷī)}{ડુંગળી (ḍungḷī)}]{\gu{ડુંગળી} (ḍungḷī) — onion}

\label{lesson:GU-C166-dungli}



\begin{quote}
Before the new one: say the Gujarati for to chop, then the Gujarati for to peel.
\end{quote}

\subsection*{You'll want to know: \gu{ડુંગળી}}
\textbf{\gu{ડુંગળી}} — \emph{ḍungḷī} — \textquotedblleft{}onion\textquotedblright{}.

\begin{quote}
\textbf{\gu{તે} \gu{ડુંગળી} \gu{સમારે} \gu{છે}.} — \emph{te ḍungḷī samāre che} — \textquotedblleft{}He chops an onion.\textquotedblright{}
\end{quote}

\begin{cousinweb}
The origin of \textbf{\gu{ડુંગળી}} (\emph{ḍungḷī}) is not settled. CDIAL lists no Gujarati form of it, and EtymDB has no entry for it.
\end{cousinweb}

\begin{grammarlens}[title={What you've built}]
One more word for everyday life at A2.
\end{grammarlens}

\subsection*{Your turn}
\begin{itemize}
\raggedright
  \item \emph{Say it:} \emph{ḍungḷī}
  \item \emph{Say it:} \emph{ḍungḷī}, once more
  \item \emph{Say it:} \emph{te ḍungḷī samāre che} — He chops an onion.
  \item \emph{Recall:} say the Gujarati for to chop, then the Gujarati for to peel, then say \emph{ḍungḷī} again
\end{itemize}

\begin{tcolorbox}[breakable,colback=teal!4,colframe=teal!35!black,title={Before you move on}]
What does \gu{ડુંગળી} mean? (\textquotedblleft{}Onion\textquotedblright{}.) Say it once more.
\end{tcolorbox}
\section[\texorpdfstring{\gu{કોથમીર} (kothmīr)}{કોથમીર (kothmīr)}]{\gu{કોથમીર} (kothmīr) — fresh coriander}

\label{lesson:GU-C166-kothmir}



\begin{quote}
Before the new one: say the Gujarati for to peel, then the Gujarati for onion.
\end{quote}

\subsection*{You'll want to know: \gu{કોથમીર}}
\textbf{\gu{કોથમીર}} — \emph{kothmīr} — \textquotedblleft{}fresh coriander\textquotedblright{}.

\begin{quote}
\textbf{\gu{કોથમીર} \gu{લીલી} \gu{અને} \gu{તાજી} \gu{છે}.} — \emph{kothmīr līlī ane tājī che} — \textquotedblleft{}The coriander is green and fresh.\textquotedblright{}
\end{quote}

\begin{cousinweb}
\textbf{\gu{કોથમીર}} (\emph{kothmīr}) is a tadbhava. It comes from Sanskrit \emph{kustumbarī}, \textquotedblleft{}coriander\textquotedblright{}, which CDIAL counts as non-Aryan in origin, and CDIAL lists the Gujarati \emph{kothmīr} (CDIAL 3380).
\end{cousinweb}

\begin{grammarlens}[title={What you've built}]
One more word for everyday life at A2.
\end{grammarlens}

\subsection*{Your turn}
\begin{itemize}
\raggedright
  \item \emph{Say it:} \emph{kothmīr}
  \item \emph{Say it:} \emph{kothmīr}, once more
  \item \emph{Say it:} \emph{kothmīr līlī ane tājī che} — The coriander is green and fresh.
  \item \emph{Recall:} say the Gujarati for to peel, then the Gujarati for onion, then say \emph{kothmīr} again
\end{itemize}

\begin{tcolorbox}[breakable,colback=teal!4,colframe=teal!35!black,title={Before you move on}]
What does \gu{કોથમીર} mean? (\textquotedblleft{}Fresh coriander\textquotedblright{}.) Say it once more.
\end{tcolorbox}
\section[\texorpdfstring{\gu{વાટવું} (vāṭvũ)}{વાટવું (vāṭvũ)}]{\gu{વાટવું} (vāṭvũ) — to grind, to crush}

\label{lesson:GU-C166-vatvu}



\begin{quote}
Before the new one: say the Gujarati for onion, then the Gujarati for fresh coriander.
\end{quote}

\subsection*{You'll want to know: \gu{વાટવું}}
\textbf{\gu{વાટવું}} — \emph{vāṭvũ} — \textquotedblleft{}to grind, to crush\textquotedblright{}.

\begin{quote}
\textbf{\gu{તે} \gu{આદુ} \gu{અને} \gu{મરચું} \gu{વાટે} \gu{છે}.} — \emph{te ādu ane marcũ vāṭe che} — \textquotedblleft{}She grinds ginger and chilli.\textquotedblright{}
\end{quote}

\begin{cousinweb}
\textbf{\gu{વાટવું}} (\emph{vāṭvũ}) is a tadbhava. It comes from Sanskrit \emph{vartayati}, \textquotedblleft{}causes to turn, whirls\textquotedblright{}, and CDIAL glosses the Gujarati verb \textquotedblleft{}to pound by rolling\textquotedblright{} (CDIAL 11356).
\end{cousinweb}

\begin{grammarlens}[title={What you've built}]
One more word for everyday life at A2.
\end{grammarlens}

\subsection*{Your turn}
\begin{itemize}
\raggedright
  \item \emph{Say it:} \emph{vāṭvũ}
  \item \emph{Say it:} \emph{vāṭvũ}, once more
  \item \emph{Say it:} \emph{te ādu ane marcũ vāṭe che} — She grinds ginger and chilli.
  \item \emph{Recall:} say the Gujarati for onion, then the Gujarati for fresh coriander, then say \emph{vāṭvũ} again
\end{itemize}

\begin{tcolorbox}[breakable,colback=teal!4,colframe=teal!35!black,title={Before you move on}]
What does \gu{વાટવું} mean? (\textquotedblleft{}To grind, to crush\textquotedblright{}.) Say it once more.
\end{tcolorbox}
